% ---------------------------------------------------------------------------
% NasAge: supplementary figures S5 and S6, with legends.
%
% Companion to figures_main.tex, which contains Figures 5 and 6 and does not
% reference this document. This document does reference the main figures.
%
% Compile with:  tectonic -X compile figures_supplementary.tex
%            or: pdflatex figures_supplementary.tex
% The figure PDFs must sit beside this file.
% ---------------------------------------------------------------------------
\documentclass[11pt,a4paper]{article}

\usepackage[margin=20mm]{geometry}
\usepackage{graphicx}
\usepackage{amsmath}
\usepackage{microtype}
\usepackage[hidelinks]{hyperref}

\setlength{\parindent}{0pt}
\setlength{\parskip}{6pt}

\newcommand{\figpage}[1]{%
  \clearpage
  \thispagestyle{empty}%
  \vspace*{\fill}
  \begin{center}
    \includegraphics[width=\textwidth,height=\textheight,keepaspectratio]{#1}
  \end{center}
  \vspace*{\fill}
}

\newcommand{\legend}[3]{%
  \clearpage
  {\large\textbf{Figure #1}\par}
  \vspace{2pt}
  {\textbf{#2}\par}
  \vspace{6pt}
  \small #3\par
}

\title{NasAge: Supplementary Figures S5 and S6}
\author{}
\date{Generated 25 November 2026}

\begin{document}
\maketitle

\section*{Purpose of this document}

These two figures carry the boundary conditions and the failures behind main Figures 5 and 6.
They exist so that neither main figure can be read as a complete account of what was tested.

Supplementary Figure S5 closes off three readings of Figure 5 that would otherwise remain
available: that the ageing signal is a continuous adult gradient, that it reflects cellular
senescence, and that an absent SASP signal is a biological null. Supplementary Figure S6
collects everything the mucosal programme panel of Figure 6 failed, and identifies the single
cause behind most of those failures.

Methods are as described with the main figures: every panel is drawn from a results file
produced by an analysis plan committed to version control before the corresponding data were
accessed, and no panel recomputes a model statistic.

% ===========================================================================
\figpage{figure5_supp.pdf}

\legend{S5}{Boundary conditions on the mucociliary ageing result.}{%
\textbf{(a)} The interferon score shows no gradient with age \emph{within} older adults, in
two independent cohorts ($\rho=-0.07$, $p=0.46$, $n=120$; $\rho=-0.16$, $p=0.32$, $n=39$).
With Figure 5d this localises the signal to the childhood-to-adult transition and excludes a
continuous adult decline.
\par
\textbf{(b)} Detected genes per cell in CDKN2A$^{+}$ versus CDKN2A$^{-}$ cells. p16$^{+}$
status tracks sequencing depth almost entirely: medians differ by 1{,}377 genes per cell and
181 of 187 donors show the same direction. Senescence calls from droplet data are therefore
confounded with capture efficiency, which is why the corresponding cohort is withdrawn from
biological claims.
\par
\textbf{(c)} Fraction of each module's genes detected in at least 5\% of cells, in four
cohorts and two chemistries. The dashed line is the pre-specified 60\% measurability gate.
M12 SASP falls at or near zero everywhere and is unmeasurable in droplet data; it is retired
from primary analysis rather than reported as a null result.
\par
\textbf{(d)} The three outcome-anchored tests with their pre-registered directions, verdicts
and the reason for each verdict. Predictions, direction and multiplicity were fixed by a
committed plan before the data were accessed.
\par
\textbf{Interpretation.} These four panels are boundary conditions, and each one closes off a
reading of Figure 5 that would otherwise be available. Panel (a) removes the continuous
ageing interpretation. Panel (b) removes the senescence interpretation: p16 status in droplet
data is very largely a statement about how deeply a cell was sequenced, so any biology built
on it is confounded at the measurement level. Panel (c) removes the possibility of reporting
SASP as a negative finding, because a module detected in essentially no cells has not been
tested at all, and a null from an unmeasurable module is uninformative rather than evidence
of absence. Panel (d) records what was promised in advance, so the verdicts in the main
figure can be checked against the plan rather than against the outcome.
\par
\textbf{Conclusion.} The ageing signal in Figure 5 is real but narrow: it is a
childhood-to-adult contrast in measurable modules, not a senescence phenotype and not an
adult gradient. Claims about cellular senescence from this data type require a depth-matched
design, and claims about SASP require an assay with a lower detection floor.}

% ===========================================================================
\figpage{figure6_supp.pdf}

\legend{S6}{Everything the programme panel failed.}{%
This figure exists so that Figure 6 is not read as a complete account.
\par
\textbf{(a)} Evidence matrix: for each of five candidate axes, whether it is measurable,
cell-type specific, changes with age, relates to infection, and is engaged by a vaccine.
Green is a pass, small gold is partial, orange is a failure, a dash is not tested. P5's
infection entry is \emph{partial}, not a pass, because it clears its matched null in one
cohort of three.
\par
\textbf{(b)} Declared axes versus the effective number of independent tests (Li and Ji, 2005)
in each cell type. Five declared axes behave as three to four, because PC1 carries 55--69\%
of their variance, so Benjamini--Hochberg over the declared count is the wrong correction in
both directions.
\par
\textbf{(c)} Age slopes for all five axes in all three cell types; a green marker denotes the
combinations that clear the expression-matched null. Seven of eleven tested combinations
reach $\mathrm{BH}<0.05$ yet fail that null, a general age drift that most gene sets
partially capture. Only antigen presentation is robust.
\par
\textbf{(d)} Paired change in each programme after live attenuated intranasal influenza
vaccine. The M05 interferon positive control rises as expected, confirming the assay works on
these samples, while every mucosal programme is null. Three independent designs returned this
result. \textbf{This is the panel that prevents the figure set being titled as a validated
vaccine-readiness index.}
\par
\textbf{(e)} Median pairwise correlation among the five modules that a revised panel (v2)
collapsed into one interferon axis, in the two cohorts the collapse was derived from and in
three held-out cohorts. Black rules mark the 95th percentile of random five-module subsets of
the same panel. The axis is 0.72--0.93 where it was derived and 0.34--0.50 where it was not,
beating chance in one held-out cohort of three. By the pre-registered rule the collapse does
not replicate, and the original panel remains the panel of record.
\par
\textbf{(f)} Empirical $p$ for every programme against its expression- and size-matched null,
in each cohort, under each of four scoring methods (mean $z$, median $z$, rank, PC1 of set).
Boxed cells clear the null. Specificity is scorer-dependent for most programmes, which is why
the reported figure is the worst scorer rather than the best.
\par
\textbf{(g)} Fraction of \emph{random} expression-matched gene sets that reach model $p<0.05$,
against the nominal 0.05. The bulk challenge cohorts are close to nominal (0.02--0.06); the
single-cell pseudobulk path is not (0.33 basal, 0.43 secretory, 0.76 ciliated). Where this
fraction is high a nominal $p$ from any gene set is anticonservative and the matched null is
the only valid test. This is the dominant-axis problem measured directly, and it is why the
trajectories in (c) are judged on their matched nulls rather than on their $p$-values.
\par
\textbf{Interpretation.} The failures collected here are not independent; most of them have
one cause, and panel (g) names it. A single dominant axis of expression explains 55--69\% of
the variance across these modules, so almost any gene set partially tracks it. That is why
five declared axes behave as three (b), why seven of eleven age associations are significant
yet unspecific (c), and why collapsing five modules into a named interferon axis fails out of
sample (e): those modules covary because nearly everything covaries, not because they share a
regulatory programme. The severity is wildly uneven between data types. In the bulk challenge
cohorts the test is close to nominal, which is what licenses the pneumococcal result in
Figure 6f. In the single-cell pseudobulk path it is not, reaching 0.76 in ciliated cells,
where a nominal $p$ carries almost no information and Benjamini--Hochberg over such
$p$-values is close to meaningless. Panel (d) stands apart from this: the vaccine null is not
a specificity failure but a straightforward negative result, and the positive control
establishes that it is a real null rather than a failed assay.
\par
\textbf{Conclusion.} Gene-set age associations in nasal single-cell data must be validated
against expression-matched nulls before they are believed; in these cohorts the nominal
$p$-value is not a usable guide, and most associations that pass correction do not survive
that check. Only M06 antigen presentation does. The revised panel reorganisation is
unsupported outside its derivation cohorts and the original panel remains the panel of
record. No mucosal programme is engaged by live attenuated influenza vaccination.}

\end{document}
